% GENERATED FILE. Edit canonical lessons, then run npm run generate:books.
% canonical-source-hash: fnv1a64:11b9a4340c90baca
% canonical-lessons: JA-C491-suiji, JA-C491-jaguchi, JA-C491-konro, JA-C491-suidousui, JA-C491-futtou-suru, JA-R491-review-486-491

\chapter{Water and Fire in the Kitchen}
\label{ch:ja-water-and-fire-in-the-kitchen}
\hlchaptermodality{491}

\begin{chapteropening}
\textbf{By the end of this chapter:} \emph{I can talk about water and heat in the kitchen.}

Complete the last lesson of chapter 491: I can talk about water and heat in the kitchen.
\end{chapteropening}

\section[\texorpdfstring{\ja{すいじ} (suiji)}{すいじ (suiji)}]{\ja{すいじ} (suiji) — kitchen work; cooking}

\label{lesson:JA-C491-suiji}



\begin{quote}
Before the new one: say the Japanese for animal feed, then the Japanese for to jump.
\end{quote}

\subsection*{You'll want to know: \ja{すいじ}}
\textbf{\ja{すいじ}} — \emph{suiji} — \textquotedblleft{}kitchen work; cooking\textquotedblright{}.

\begin{quote}
\textbf{\ja{ちちは} \ja{にちようびに} \ja{すいじを} \ja{します。}} — \emph{chichi wa nichiyōbi ni suiji o shimasu} — \textquotedblleft{}My father does the cooking on Sunday.\textquotedblright{}
\end{quote}

\begin{cousinweb}
\textbf{\ja{すいじ}} (\emph{suiji}) is Sino-Japanese (kango): two Chinese characters read with their on'yomi, readings borrowed from Chinese. They are \emph{sui}, \textquotedblleft{}cook\textquotedblright{} (Middle Chinese \emph{tsyhwe}) and \emph{ji}, \textquotedblleft{}matter\textquotedblright{} (\emph{dzriH}) (KANJIDIC2; Middle Chinese in Baxter's transcription).
\end{cousinweb}

\begin{grammarlens}[title={What you've built}]
One more word for describing what you lived and felt.
\end{grammarlens}

\subsection*{Your turn}
\begin{itemize}
\raggedright
  \item \emph{Say it:} \emph{suiji}
  \item \emph{Say it:} \emph{suiji}, once more
  \item \emph{Say it:} \emph{chichi wa nichiyōbi ni suiji o shimasu} — My father does the cooking on Sunday.
  \item \emph{Recall:} say the Japanese for animal feed, then the Japanese for to jump, then say \emph{suiji} again
\end{itemize}

\begin{tcolorbox}[breakable,colback=teal!4,colframe=teal!35!black,title={Before you move on}]
What does \ja{すいじ} mean? (\textquotedblleft{}Kitchen work; cooking\textquotedblright{}.) Say it once more.
\end{tcolorbox}
\section[\texorpdfstring{\ja{じゃぐち} (jaguchi)}{じゃぐち (jaguchi)}]{\ja{じゃぐち} (jaguchi) — a tap; a faucet}

\label{lesson:JA-C491-jaguchi}



\begin{quote}
Before the new one: say the Japanese for to jump, then the Japanese for kitchen work.
\end{quote}

\subsection*{You'll want to know: \ja{じゃぐち}}
\textbf{\ja{じゃぐち}} — \emph{jaguchi} — \textquotedblleft{}a tap; a faucet\textquotedblright{}.

\begin{quote}
\textbf{\ja{じゃぐちで} \ja{てを} \ja{あらいます。}} — \emph{jaguchi de te o araimasu} — \textquotedblleft{}I wash my hands at the tap.\textquotedblright{}
\end{quote}

\begin{cousinweb}
\textbf{\ja{じゃぐち}} (\emph{jaguchi}) mixes the two layers, native and Sino-Japanese: Sino-Japanese \emph{ja}, \textquotedblleft{}snake\textquotedblright{} (Middle Chinese \emph{zyae}), borrowed from Chinese and native \emph{guchi}, \textquotedblleft{}mouth\textquotedblright{} (JMdict, KANJIDIC2; Middle Chinese in Baxter's transcription).
\end{cousinweb}

\begin{grammarlens}[title={What you've built}]
One more word for describing what you lived and felt.
\end{grammarlens}

\subsection*{Your turn}
\begin{itemize}
\raggedright
  \item \emph{Say it:} \emph{jaguchi}
  \item \emph{Say it:} \emph{jaguchi}, once more
  \item \emph{Say it:} \emph{jaguchi de te o araimasu} — I wash my hands at the tap.
  \item \emph{Recall:} say the Japanese for to jump, then the Japanese for kitchen work, then say \emph{jaguchi} again
\end{itemize}

\begin{tcolorbox}[breakable,colback=teal!4,colframe=teal!35!black,title={Before you move on}]
What does \ja{じゃぐち} mean? (\textquotedblleft{}A tap; a faucet\textquotedblright{}.) Say it once more.
\end{tcolorbox}
\section[\texorpdfstring{\ja{こんろ} (konro)}{こんろ (konro)}]{\ja{こんろ} (konro) — a cooking stove}

\label{lesson:JA-C491-konro}



\begin{quote}
Before the new one: say the Japanese for kitchen work, then the Japanese for a tap.
\end{quote}

\subsection*{You'll want to know: \ja{こんろ}}
\textbf{\ja{こんろ}} — \emph{konro} — \textquotedblleft{}a cooking stove\textquotedblright{}.

\begin{quote}
\textbf{\ja{こんろで} \ja{ゆを} \ja{わかします。}} — \emph{konro de yu o wakashimasu} — \textquotedblleft{}I boil water on the stove.\textquotedblright{}
\end{quote}

\begin{cousinweb}
\textbf{\ja{こんろ}} (\emph{konro}) is Sino-Japanese (kango): two Chinese characters read with their on'yomi, readings borrowed from Chinese. They are \emph{kon}, \textquotedblleft{}shine\textquotedblright{} and \emph{ro}, \textquotedblleft{}hearth\textquotedblright{} (Middle Chinese \emph{lu}) (KANJIDIC2; Middle Chinese in Baxter's transcription).
\end{cousinweb}

\begin{grammarlens}[title={What you've built}]
One more word for describing what you lived and felt.
\end{grammarlens}

\subsection*{Your turn}
\begin{itemize}
\raggedright
  \item \emph{Say it:} \emph{konro}
  \item \emph{Say it:} \emph{konro}, once more
  \item \emph{Say it:} \emph{konro de yu o wakashimasu} — I boil water on the stove.
  \item \emph{Recall:} say the Japanese for kitchen work, then the Japanese for a tap, then say \emph{konro} again
\end{itemize}

\begin{tcolorbox}[breakable,colback=teal!4,colframe=teal!35!black,title={Before you move on}]
What does \ja{こんろ} mean? (\textquotedblleft{}A cooking stove\textquotedblright{}.) Say it once more.
\end{tcolorbox}
\section[\texorpdfstring{\ja{すいどうすい} (suidōsui)}{すいどうすい (suidōsui)}]{\ja{すいどうすい} (suidōsui) — tap water}

\label{lesson:JA-C491-suidousui}



\begin{quote}
Before the new one: say the Japanese for a tap, then the Japanese for a cooking stove.
\end{quote}

\subsection*{You'll want to know: \ja{すいどうすい}}
\textbf{\ja{すいどうすい}} — \emph{suidōsui} — \textquotedblleft{}tap water\textquotedblright{}.

\begin{quote}
\textbf{\ja{この} \ja{まちの} \ja{すいどうすいは} \ja{おいしいです。}} — \emph{kono machi no suidōsui wa oishii desu} — \textquotedblleft{}The tap water in this town tastes good.\textquotedblright{}
\end{quote}

\begin{cousinweb}
\textbf{\ja{すいどうすい}} (\emph{suidōsui}) is Sino-Japanese (kango): three Chinese characters read with their on'yomi, readings borrowed from Chinese. They are \emph{sui}, \textquotedblleft{}water\textquotedblright{} (Middle Chinese \emph{sywijX}), \emph{dō}, \textquotedblleft{}road\textquotedblright{} (\emph{dawX}) and \emph{sui}, \textquotedblleft{}water\textquotedblright{}, again: the water of the \emph{suidō}, the water main (KANJIDIC2; Middle Chinese in Baxter's transcription).
\end{cousinweb}

\begin{grammarlens}[title={What you've built}]
One more word for describing what you lived and felt.
\end{grammarlens}

\subsection*{Your turn}
\begin{itemize}
\raggedright
  \item \emph{Say it:} \emph{suidōsui}
  \item \emph{Say it:} \emph{suidōsui}, once more
  \item \emph{Say it:} \emph{kono machi no suidōsui wa oishii desu} — The tap water in this town tastes good.
  \item \emph{Recall:} say the Japanese for a tap, then the Japanese for a cooking stove, then say \emph{suidōsui} again
\end{itemize}

\begin{tcolorbox}[breakable,colback=teal!4,colframe=teal!35!black,title={Before you move on}]
What does \ja{すいどうすい} mean? (\textquotedblleft{}Tap water\textquotedblright{}.) Say it once more.
\end{tcolorbox}
\section[\texorpdfstring{\ja{ふっとうする} (futtō suru)}{ふっとうする (futtō suru)}]{\ja{ふっとうする} (futtō suru) — to come to the boil}

\label{lesson:JA-C491-futtou-suru}



\begin{quote}
Before the new one: say the Japanese for a cooking stove, then the Japanese for tap water.
\end{quote}

\subsection*{You'll want to know: \ja{ふっとうする}}
\textbf{\ja{ふっとうする}} — \emph{futtō suru} — \textquotedblleft{}to come to the boil\textquotedblright{}.

\begin{quote}
\textbf{\ja{やかんの} \ja{ゆが} \ja{ふっとうしたら、たまごを} \ja{いれて} \ja{ください。}} — \emph{yakan no yu ga futtōshitara, tamago o irete kudasai} — \textquotedblleft{}When the water boils, put the eggs in.\textquotedblright{}
\end{quote}

Polite \textbf{\ja{ふっとうします}}, \emph{futtōshimasu}; a \ja{する} verb. Past (polite): \textbf{\ja{やかんの} \ja{みずが} \ja{ふっとうしました。}}, \emph{yakan no mizu ga futtōshimashita}.

\begin{cousinweb}
\textbf{\ja{ふっとうする}} (\emph{futtō suru}) is the Sino-Japanese (kango) noun \emph{futtō} with the native verb \emph{suru}, \textquotedblleft{}to do\textquotedblright{}: two Chinese characters read with their on'yomi, readings borrowed from Chinese. They are \emph{fu}, \textquotedblleft{}seethe\textquotedblright{} (Middle Chinese \emph{pj+jH}) and \emph{tō}, \textquotedblleft{}leap\textquotedblright{} (\emph{dong}) (KANJIDIC2; Middle Chinese in Baxter's transcription).
\end{cousinweb}

\begin{grammarlens}[title={What you've built}]
One more word for describing what you lived and felt.
\end{grammarlens}

\subsection*{Your turn}
\begin{itemize}
\raggedright
  \item \emph{Say it:} \emph{futtō suru}
  \item \emph{Say it:} \emph{futtō suru}, once more
  \item \emph{Say it:} \emph{yakan no yu ga futtōshitara, tamago o irete kudasai} — When the water boils, put the eggs in.
  \item \emph{Recall:} say the Japanese for a cooking stove, then the Japanese for tap water, then say \emph{futtō suru} again
\end{itemize}

\begin{tcolorbox}[breakable,colback=teal!4,colframe=teal!35!black,title={Before you move on}]
What does \ja{ふっとうする} mean? (\textquotedblleft{}To come to the boil\textquotedblright{}.) Say it once more.
\end{tcolorbox}
\section[(dialogue)]{Looking back: from Parts and Fittings to Water and Fire in the Kitchen}

\label{lesson:JA-R491-review-486-491}



\begin{quote}
Say the Japanese for a part (of a machine) and a metal fitting.
\end{quote}

\subsection*{Your turn}
Say these aloud:

\begin{itemize}
\raggedright
  \item a part (of a machine), a metal fitting, to rust, to fit in, to peel off
  \item business, a merchant, a shop manager, sales, an employee
  \item architecture, timber, a brick, iron, construction work
  \item farming, a harvest, to plough, to ripen, wheat
  \item a fisherman, the fishing industry, a fishing boat, animal feed, to jump
  \item kitchen work, a tap, a cooking stove, tap water, to come to the boil
\end{itemize}

\begin{tcolorbox}[breakable,colback=teal!4,colframe=teal!35!black,title={Before you move on}]
A part (of a machine)? (\textbf{\ja{ぶひん}}, \emph{buhin}.) A metal fitting? (\textbf{\ja{かなぐ}}, \emph{kanagu}.) To rust? (\textbf{\ja{さびる}}, \emph{sabiru}.) To fit in? (\textbf{\ja{はめる}}, \emph{hameru}.) To peel off? (\textbf{\ja{はがす}}, \emph{hagasu}.) Business? (\textbf{\ja{しょうばい}}, \emph{shōbai}.) A merchant? (\textbf{\ja{しょうにん}}, \emph{shōnin}.) A shop manager? (\textbf{\ja{てんちょう}}, \emph{tenchō}.) Sales? (\textbf{\ja{うりあげ}}, \emph{uriage}.) An employee? (\textbf{\ja{じゅうぎょういん}}, \emph{jūgyōin}.) Architecture? (\textbf{\ja{けんちく}}, \emph{kenchiku}.) Timber? (\textbf{\ja{ざいもく}}, \emph{zaimoku}.) A brick? (\textbf{\ja{れんが}}, \emph{renga}.) Iron? (\textbf{\ja{てつ}}, \emph{tetsu}.) Construction work? (\textbf{\ja{こうじ}}, \emph{kōji}.) Farming? (\textbf{\ja{のうぎょう}}, \emph{nōgyō}.) A harvest? (\textbf{\ja{しゅうかく}}, \emph{shūkaku}.) To plough? (\textbf{\ja{たがやす}}, \emph{tagayasu}.) To ripen? (\textbf{\ja{みのる}}, \emph{minoru}.) Wheat? (\textbf{\ja{こむぎ}}, \emph{komugi}.) A fisherman? (\textbf{\ja{りょうし}}, \emph{ryōshi}.) The fishing industry? (\textbf{\ja{ぎょぎょう}}, \emph{gyogyō}.) A fishing boat? (\textbf{\ja{ぎょせん}}, \emph{gyosen}.) Animal feed? (\textbf{\ja{えさ}}, \emph{esa}.) To jump? (\textbf{\ja{はねる}}, \emph{haneru}.) Kitchen work? (\textbf{\ja{すいじ}}, \emph{suiji}.) A tap? (\textbf{\ja{じゃぐち}}, \emph{jaguchi}.) A cooking stove? (\textbf{\ja{こんろ}}, \emph{konro}.) Tap water? (\textbf{\ja{すいどうすい}}, \emph{suidōsui}.) To come to the boil? (\textbf{\ja{ふっとうする}}, \emph{futtō suru}.)
\end{tcolorbox}
